\section{本章のまとめ}
\label{sec:hhl-chapter-summary}

本章では，HHLアルゴリズムを用いた量子線形回帰に必要な理論と処理手順を整理した．

線形回帰の最小二乗解は，$A=X^{\mathsf{T}}X$，$\bm{b}=X^{\mathsf{T}}\bm{y}$とおくことで，
線形方程式$A\bm{w}=\bm{b}$の解として表される．
$X$が最大列階数を持つ場合，$A$は実対称正定値行列となり，
HHLアルゴリズムが対象とするエルミートかつ可逆な行列の条件を満たす．
可逆性や数値安定性に問題がある場合は，リッジ回帰による正則化を利用できる．

HHLアルゴリズムでは，量子位相推定によって$A$の固有値を推定し，
制御回転によって固有値の逆数を補助量子ビットの振幅へ反映する．
その後，逆量子位相推定と補助量子ビットの測定を行うことで，
$A^{-1}\ket{b}$に比例する解状態を得る．
線形回帰へ適用した場合，この解状態は規格化された回帰係数状態$\ket{w}$となる．

新しい入力状態$\ket{x_{*}}$と$\ket{w}$の内積を評価すれば，
回帰係数の全成分を読み出さずに予測値を求められる．
一方，HHLアルゴリズムの計算上の利点には，行列の疎性，良好な条件数，効率的な状態準備，
ハミルトニアンシミュレーション，および少数の期待値だけを必要とする出力形式が前提となる．
特に，正規方程式によって条件数が二乗されることと，$X^{\mathsf{T}}X$が密行列になり得ることは，
量子線形回帰を設計するうえで重要な課題である．

以上により，線形回帰問題の準備からHHLアルゴリズムの実行，回帰予測の評価までの流れが得られた．
今後，具体的なデータと小規模な行列を用いて各処理を実装する際には，
本章で示した前提条件と誤差要因を確認しながら検証する必要がある．
